\markdownRendererDocumentBegin
\markdownRendererSectionBegin
\markdownRendererHeadingOne{Identificar dos arboles isomorficos}\markdownRendererInterblockSeparator
{}Determinar si dos árboles son isomorfos implica verificar si existe una biyección que preserve adyacencias. El algoritmo \markdownRendererDollarSign{}AHU (Aho, Hopcroft, Ullman)\markdownRendererDollarSign{} resuelve esto en \markdownRendererDollarSign{}O(n)\markdownRendererDollarSign{} mediante hashing canónico: asigna etiquetas únicas a subárboles desde las hojas hacia la raíz, creando una representación invariante bajo isomorfismo. Dos árboles son isomorfos si sus representaciones canónicas coinciden, siendo útil en detección de patrones estructurales.\markdownRendererInterblockSeparator
{}\markdownRendererInputFencedCode{./_markdown_main/2f816965fbd6325825132dedc543b12e.verbatim}{cpp}{cpp}
\markdownRendererSectionEnd \markdownRendererDocumentEnd\markdownRendererDocumentBegin
\markdownRendererSectionBegin
\markdownRendererHeadingOne{Identificar dos arboles isomorficos}\markdownRendererInterblockSeparator
{}Determinar si dos árboles son isomorfos implica verificar si existe una biyección que preserve adyacencias. El algoritmo \markdownRendererDollarSign{}AHU (Aho, Hopcroft, Ullman)\markdownRendererDollarSign{} resuelve esto en \markdownRendererDollarSign{}O(n)\markdownRendererDollarSign{} mediante hashing canónico: asigna etiquetas únicas a subárboles desde las hojas hacia la raíz, creando una representación invariante bajo isomorfismo. Dos árboles son isomorfos si sus representaciones canónicas coinciden, siendo útil en detección de patrones estructurales.\markdownRendererInterblockSeparator
{}\markdownRendererInputFencedCode{./_markdown_main/2f816965fbd6325825132dedc543b12e.verbatim}{cpp}{cpp}
\markdownRendererSectionEnd \markdownRendererDocumentEnd